% Part: first-order-logic
% Chapter: syntax
% Section: intro-syntax

\documentclass[../../../include/open-logic-section]{subfiles}

\begin{document}

\olfileid[hi]{fol}{syn}{itx}

\olsection{परिचय}

प्रथम-क्रम तर्क के सिद्धांत और अधिसिद्धांत को विकसित करने के लिए हमें पहले
उसकी अभिव्यक्तियों के वाक्यविन्यास और अर्थविज्ञान को परिभाषित करना होगा।
प्रथम-क्रम तर्क की अभिव्यक्तियाँ पद और !!{formula}s हैं। पद !!{variable}s,
!!{constant}s और !!{function}s से बनते हैं। इसके बाद !!^{formula}s,
!!{predicate}s और पदों से बनते हैं (इनसे सबसे छोटे, ``परमाण्विक''
!!{formula}s बनते हैं); फिर तार्किक संयोजकों और परिमाणकों का उपयोग करके
परमाण्विक !!{formula}s से अधिक जटिल सूत्र बनाए जा सकते हैं। निर्माण-नियम
देने के अनेक ढंग हैं; यहाँ हम केवल एक संभव ढंग देते हैं। अन्य पद्धतियाँ भिन्न
प्रतीक चुनेंगी, संयोजकों के भिन्न समुच्चयों को आदिम मानेंगी और कोष्ठकों का
भिन्न प्रकार से उपयोग करेंगी (या तथाकथित पोलिश संकेतन की तरह उनका बिल्कुल
उपयोग नहीं करेंगी)। फिर भी सभी दृष्टिकोणों में यह बात समान है कि निर्माण-नियम
पदों और !!{formula}s के समुच्चय को \emph{आगमनात्मक रूप से} परिभाषित करते हैं।
परिभाषा ठीक ढंग से दी गई हो तो निर्माण-नियमों के अनुसार प्रत्येक अभिव्यक्ति
मूलतः केवल एक ही प्रकार से बन सकती है। अभिव्यक्तियों को \emph{अद्वितीय रूप से
पठनीय} बनाने वाली आगमनात्मक परिभाषा के कारण हम उसी विधि---आगमनात्मक
परिभाषा---से इन अभिव्यक्तियों को अर्थ दे सकते हैं।

\end{document}
